% Proposed replacement fragment, not inserted into the open response document.
% Requires xcolor, natbib and AE1_references.bib.
{\color{blue}
\noindent\textbf{Response:}
We thank the Associate Editor for suggesting this perspective. We agree that the distinction between reference knowledge and category-specific procedural expertise clarifies the positioning of LEAN-LLM-OPT.

Prior work distinguishes declarative knowledge, such as facts relevant to a problem, from procedural knowledge, such as strategies for performing a task \citep{ae1_li2024metacognitive}. In LEAN-LLM-OPT, Ref-Data provides external, case-based modeling knowledge: reference problem descriptions, associated data, expert formulations, and annotations identifying relevant data. These cases illustrate model structure and the correspondence between data and model components. They can also convey procedural information when incorporated into a worked demonstration; we therefore do not regard reference knowledge as exclusively declarative.

Retrieval-augmented generation (RAG) describes how external content is retrieved and used to condition generation \citep{ae1_lewis2020rag}. OpenAI's documentation likewise distinguishes this mechanism from Agent Skills: RAG supplies relevant context, whereas a skill packages reusable instructions and supporting resources for carrying out a task \citep{ae1_openai2026accuracy,ae1_openai2026skills}. These concepts concern different aspects of a system and can be combined. A skill can prescribe how to search for, interpret, and apply retrieved material; OpenAI explicitly illustrates skills that combine search and fetch tools into a research workflow \citep{ae1_openai2026skilltools}. Thus, a skill may organize a RAG workflow, but retrieving reference content does not by itself constitute a modeling skill.

In LEAN-LLM-OPT, reference retrieval selects relevant cases, while category-conditioned instructions specify how the agent should use them together with the current problem and its data. For the predefined categories, selected cases are assembled into demonstrations of data access and formulation. The model-generation agent then follows this guidance to obtain instance data and construct the target model. For example, the NRM route combines a reference demonstration with instructions for obtaining data, defining a symbolic formulation, and specifying exact data-to-model mappings. The general route instead constructs an instance-specific abstract modeling plan from the current query and data schema. Reference cases provide modeling examples; target datasets provide the parameters of the current instance.

This reusable procedural component is \emph{skill-like}, consistent with the characterization of skills as procedural artifacts with instructions, supporting resources, and applicability conditions \citep{ae1_zhou2026skillsurvey}. Reuse concerns the modeling method, not an identical action sequence or output across instances. Our implementation embeds this method in prompts, workflow-construction code, and data tools rather than standalone \texttt{SKILL.md} packages. Separately, it does not acquire or revise a persistent skill library from completed trajectories. This distinguishes its design-time curation from the workflow induction studied in Agent Workflow Memory and the archetype-based skill distillation and refinement studied in OptSkills \citep{ae1_wang2025awm,ae1_yang2026optskills}.

Accordingly, we position LEAN-LLM-OPT as a \emph{reference-guided agentic workflow construction framework} that combines case-based modeling knowledge with reusable, category-conditioned procedural guidance for optimization modeling with external data. The agents apply these resources through reasoning and tool use; the resources could also be used with a fine-tuned model. Our contribution concerns their organization and application at inference time, rather than establishing that references replace agents or are universally preferable to fine-tuning.

We have added \emph{Knowledge and Skills in LLM Agents} to Section~1.2 (\emph{Related Works}) to explain knowledge, RAG, and skills and their connections to recent literature. We have also added \emph{Connecting Reference Knowledge and Modeling Skills} to Section~3.3 (\emph{Novelties in Our Design and Key Messages}) to explain their roles in LEAN-LLM-OPT and distinguish procedural reuse from standalone skill packaging and trajectory-based skill learning.
}
